\documentclass[../main.tex]{subfiles}
\begin{document}
%%%%%%%%%%%%%%%%%%%%%%%%%%%%%%%%%%%%%%%%%%%%%%%%%%%%%%%%%%%%%%%%%%%%%%%%%%%%%%%%%%%%%%%%%%%%%%%%%%%
\chapter{Working With the Inspectors}

Along the right side of the window, you'll find the \textbf{Inspectors} sidebar. You use the
Inspectors to change the style attributes for the objects on the canvas, set the properties of
the canvas itself, and to include document-specific metadata and define how the OmniGraffle
file is saved.

\begin{ogtip}
  If you don't see the Inspector sidebar, click \ogicon{12.5pt}{res/og7_button_inspectors.png} in
  the toolbar or choose \textbf{Inspectors} $\blacktriangleright$ \textbf{Hide/Show Inspectors}
  (\textbf{Shift-Command-I}).
\end{ogtip}

The Inspectors sidebar has five tabs along its top, used for categorizing the individual
\textit{inspector panes} based on their purpose:

\ogfigure{0.800}{res/og7_insp_tabs.png}

\begin{enumerate}
  \item \textbf{Object} inspectors are used for styling the objects on the canvas.
  \item \textbf{Properties} inspectors are used for defining the properties of the objects on
    the canvas.
  \item \textbf{Canvas} inspectors are used for defining the canvas itself, including units of
    measure and the alignment grid.
  \item \textbf{Document} inspectors are used for setting how the OmniGraffle file is saved,
    setting printer margins, and for adding document metadata.
  \item \textbf{Stencil Browser} is used for finding and placing stencils (pre-styled objects)
    on the canvas.
\end{enumerate}

\begin{ognote}
  Since the Stencil Browser is not an inspector, it won't be covered in this chapter. See Using,
  Curating, and Creating Stencils for more details about the Stencil Browser and stencils in
  general.
\end{ognote}

To make it easy for you to access the inspectors, we have conveniently mapped the Inspector tabs
to the following keyboard shortcuts:

\begin{itemize}
  \item Use \textbf{Command--1} to view the Object inspectors.
  \item Use \textbf{Command--2} to view the Properties inspectors.
  \item Use \textbf{Command--3} to view the Canvas inspectors.
  \item Use \textbf{Command--4} to view the Document inspectors.
\end{itemize}

Each tab contains a series of \textit{inspector panes}, for such things as adding a fill to an
object or setting a background image for the project you're working on.

To open an inspector pane, hover over the title with the mouse, and then click on the disclosure
triangle to open and use the inspector. To open one particular inspector (for example the
\textbf{Fill} inspector) and close any other open inspectors, \textbf{Option-click} on the
disclosure triangle. To open or close all of the inspectors, \textbf{Shift-click} on one of
the disclosure triangles.

\ogfigure{0.800}{res/og7_inspectors01.png}

\begin{ogtip}
  You can configure OmniGraffle's preferences so that only the inspectors you need are displayed
  in the sidebar. Choose \textbf{OmniGraffle} $\blacktriangleright$ \textbf{Preferences}
  $\blacktriangleright$ \textbf{General} $\blacktriangleright$ \textbf{Inspectors} and turn on
  \textbf{Hide inspectors that don't apply}.
\end{ogtip}

\section{Using the Inspectors}

By default, OmniGraffle's inspectors reside in a sidebar just to the right of the Canvas.
However, you can also open the inspectors as a floating window, or as four floating palettes,
one for each tab. The different ways of viewing the Inspectors are known as Workspaces,
because they help define how you work in OmniGraffle.

You can quickly switch between Workspaces by using the following keyboard shortcuts:

\begin{itemize}
  \item \textbf{Sidebar}: F2
  \item \textbf{Floating}: F3
  \item \textbf{Palettes}: F4
\end{itemize}

\begin{ognote}
  If you are using a MacBook Pro, MacBook Air, or an
  \href{https://www.apple.com/shop/product/MLA22LL/A/magic-keyboard-us-english}{Apple Magic
  Keyboard} with your Mac, you also need to press the \textbf{fn} key---located at the
  lower-left corner of your keyboard---along with the applicable function key. If you have an
  \href{https://www.apple.com/shop/product/MQ052LL/A/magic-keyboard-with-numeric-keypad-us-english}{Apple
  Magic Keyboard with Numeric Keypad}, the \textbf{fn} key is located below the \textbf{F13}
  key.
\end{ognote}

\subsection{Interacting with the Inspectors}

In all of the inspectors, you can click in a text field that contains a number, and then press
the up or down arrow keys to increment or decrement the number.

When it comes to rotating objects, you have a couple options. First select the object and then
go to the \textbf{Geometry} inspector. Click and spin the rotation dial and watch the object
revolve around its center point. After clicking the rotation dial, keep the mouse button held
down while moving the pointer away from the control to gain more rotational precision.

\ogfigure{0.800}{res/og7_inspectors05.png}

If you have enabled the use of Multi-Touch gestures (\textbf{OmniGraffle} $\blacktriangleright$
\textbf{Preferences} $\blacktriangleright$ \textbf{Multi-Touch}), you can use a two-fingered
spin gesture to rotate an object on the canvas.

\begin{ognote}
  The use of Multi-Touch gestures requires a touch-enabled trackpad. Most Apple laptops support
  Multi-Touch, as does Apple's Magic Trackpad.
\end{ognote}

Select and touch an object with two fingers, and then rotate the fingers (or pivot one around
the other) to rotate the object. While rotating the object, you'll notice that a blue angle
indicator appears within the object, and a tiny popover appears showing you the degree of the
angle.

\ogfigure{0.800}{res/og7_inspectors06.png}

Any color swatch in the Colors window can be dragged to an object on the canvas or to an
inspector's color control. When you drag a color swatch to an object, you can drop the color
on the object's fill, stroke, or text (label) color.

\ogfigure{0.800}{res/og7_inspectors07.png}

In places where you can enter measurements, such as in the Geometry inspector, values are shown
in the current ruler units, or, if there is no unit scale, in the Units inspector. You can,
however, enter values in any unit type that is available in the ruler, such as miles or
kilometers. As soon as you finish entering the value, OmniGraffle converts it to the correct
units automatically.

\ogfigure{0.800}{res/og7_inspectors08.png}

You can also perform simple arithmetic in the input fields; for example, add (\textbf{+}),
subtract (\textbf{$-$}), multiply (\textbf{*}), or divide (\textbf{/}).

\ogfigure{0.800}{res/og7_inspectors09.png}

If you don't specify a unit type in your math (for example, \textbf{(2+2)*25}), OmniGraffle uses
the Unit Type specified in the Units inspector. You can, however, enter a different Unit Type
as part of the arithmetic (for example, \textbf{(2in+2cm)*25}) and OmniGraffle will convert
the Unit Type you have specified as part of the math to the Unit Type specified in the Units
inspector.

You can also enter unit values that are different from the \textbf{Ruler Units} specified in the
\textbf{Units} inspector. For example, if you have your project set up to use \textbf{inches}
for the Ruler Units, and you enter \textbf{2.54 cm} in a field, OmniGraffle would convert that
to \textbf{1 in}.

\section{The Object Inspectors}

Use the \textbf{Object inspectors} to view and change details about the objects on the Canvas.
You can quickly access the Object inspectors by pressing \textbf{Command-1}, or by selecting
any object on the canvas.

\subsection{Coloring Shapes with the Fill Inspector}

Use the Fill inspector to choose a color, gradient, or pattern to fill the inside of the
selected shapes. Depending on the type of fill you've selected in the \textbf{Fill Type}
pop-up menu, the Fill inspector adapts to provide the controls you need to adjust the fill's
settings.

\ogfigure{0.800}{res/og7_inspectors10.png}

\ogfigure{0.800}{res/og7_inspectors11.png}

The Fill inspector has the following controls:

\begin{itemize}
  \item \textbf{Fill Type}: Use the pop-up menu at the top of the inspector to choose from the
    nine fill types, including: Solid, Linear Gradient, Radial Gradient, Double Linear Gradient,
    Double Radial Gradient, Stipple, Marker, Squiggle, and Plastic. The \textbf{No Fill} option
    is also available as an option for objects where transparency is desired.

    With the \textbf{Fill Type} set to \textbf{Solid}, you'll also see a series of default color
    swatches that you can choose from. The larger swatches along the top range from black to
    white with varying shades of gray in between, with smaller color swatches underneath.
  \item \textbf{Fill Color}: Located underneath the Fill Type pop-up menu is the Fill Color
    pop-up. Clicking the arrows lets you choose from one of 12 preset colors, or you can click
    on the color well to the left of the arrows to open the Colors palette to blend your own
    color.
  \item \textbf{Color Swatches}: Beneath the Fill Color picker is a palette of color swatches
    that you can click and drag to quickly apply colors to objects on the canvas, or anywhere
    else that can receive a color swatch application.

    These color swatches start with a default palette set, but can be customized by dragging from
    a color well in the color picker and dropping onto a spot in the palette. Other drag sources
    for color (such as the style well in the bottom toolbar) also work for this.

    To restore the default color to a palette swatch, control-click it to open the contextual menu
    and choose \textbf{Restore Default Color}.
\end{itemize}

\begin{ogpro}
  If you have OmniGraffle Pro, the Fill inspector has two additional pop-up menus along the
  bottom: \textbf{Fill Blend Type} and \textbf{Distortion Effect}.

  \ogfigure{0.800}{res/og7_inspectors12.png}

  \begin{itemize}
    \item \textbf{Blending Mode}: The Blending Mode pop-up lets you choose how to blend the
      currently selected object's fill with objects on the same canvas. By default, this is set
      to \textbf{No Effect}.

      \ogfigure{0.800}{res/og7_inspectors13.png}
    \item \textbf{Distortion Effect Type}: The Fill inspector includes three Distortion Effect
      Types: \textbf{Blur}, \textbf{Pixelate}, and \textbf{Magnify}. For example, you could
      apply the \textbf{Magnify} effect to an object that sits atop an image so you can use it
      like a magnifying glass.

      \ogfigure{0.800}{res/og7_inspectors14.png}

      \ogfigure{0.800}{res/og7_inspectors15.png}

      \ogfigure{0.800}{res/og7_inspectors16.png}
  \end{itemize}

  You can also add a \textbf{pattern fill} to the objects in your projects. When you click on
  the color portion of a Fill Color pop-up menu, that opens the Colors window. Click the tab
  on the right, the \textbf{Pattern Palette}, to choose from a pattern to fill the objects.

  \ogfigure{0.800}{res/og7_inspectors18.png}
\end{ogpro}

Some additional notes about the Fill inspector:

\begin{itemize}
  \item The gradient rotation control becomes active when you choose one of the Linear Gradient
    fills; enter a number of degrees in the field or drag the circular control to set the angle
    of the gradient. The blend position control becomes active when you choose a radial
    gradient; drag its crosshair to move the gradient's center.
  \item If you have chosen \textbf{Linear Gradient} or \textbf{Radial Gradient}, you can click
    the swap button between the color wells to exchange the two colors.

    \ogfigure{0.800}{res/og7_inspectors17.png}
  \item Drag the \textbf{Gradient Midpoint slider} next to the color wells to adjust the
    gradient.
  \item If you choose one of the Gradient fills, two or three color wells become available.
    Click them to choose the colors to blend together.
\end{itemize}

You can also edit the fill color of a canvas itself using the \textbf{Canvas Fill} inspector.

\subsection{Defining Line and Path Behaviors with the Stroke Inspector}

Use the Stroke inspector to change the appearance of the selected connection lines or of the
line drawn around the edge of the selected shapes.

\ogfigure{0.800}{res/og7_inspectors19.png}

\begin{itemize}
  \item Use the dropdown menu to choose one of the following stroke types: No Stroke, Single
    Stroke, Double Stroke, Freehand Stroke, Plastic Stroke, Inner Stroke, or Outer Stroke.
  \item Click the color well to choose a color for the stroke.
  \item Enter a number in the Thickness field or use the keyboard arrows to choose how thick the
    stroke should be.
  \item With any Stroke Type except \textbf{Plastic Stroke} selected, you'll find three sets of
    controls for further defining the stroke:

    \ogfigure{0.800}{res/og7_inspectors20.png}

    \begin{itemize}
      \item Use the pop-up menu on the left to select a \textbf{Stroke Pattern}, ranging from
        solid, to dashed, to dotted, and everything in between. There are 25 stroke patterns to
        choose from.
      \item The center set of controls determines how the stroke appears at its corners:
        \textbf{Miter} creates a sharp corner, \textbf{Round} creates a soft corner, and
        \textbf{Bevel} creates a cut-off corner.
      \item The controls on the right determine how the stroke appears at its ends: The
        \textbf{Butt} option ends the line by cutting across its endpoint at an angle
        perpendicular to the line itself, while the \textbf{Round} and \textbf{Square} options let
        the line extend past its endpoint based on the stroke width.
    \end{itemize}
  \item With the \textbf{Plastic Stroke} type selected, the Stroke inspector takes on a
    different set of controls:

    \ogfigure{0.800}{res/og7_inspectors21.png}

    The middle row of controls lets you set a color for the Plastic Stroke's \textbf{Stroke} and
    \textbf{Highlight} colors. Use the \textbf{Highlight angle} wheel on the left to change the
    ``lighting'' angle for the highlight. Use the slider on the right to make the Plastic Stroke
    flatter or to give it more height.
\end{itemize}

\subsection{Casting Shadows with the Shadow Inspector}

Use the Shadow inspector to drop a shadow behind the selected objects.

\ogfigure{0.800}{res/og7_inspectors22.png}

\begin{itemize}
  \item Use the pop-up menu to choose whether the selected object has a shadow, and whether the
    shadow appears in front of or behind objects on the same layer, or if the object takes on an
    inner shadow. If you choose to let shadows appear in front of objects on the same layer, the
    ordering of objects matters; you can reorder objects with the \textbf{Bring} and
    \textbf{Send} commands in the \textbf{Arrange} menu.
  \item Enter a number in the blur field or use the keyboard arrows to choose how precise or
    blurry the shadow should be.
  \item Click the color well to choose a color for the shadow; partially transparent black tends
    to work best.
  \item Drag the crosshair in the offset control or enter numbers in the offset fields to set
    how far away from the object the shadow should fall.
\end{itemize}

\subsection{Choosing a Shape with the Shape Inspector}

With a shape object selected, use this inspector to transform it in various ways.

\ogfigure{0.800}{res/og7_inspectors23.png}

\begin{itemize}
  \item To change the shape (for example, from a rectangle to a circle), click the menu at the
    top of the inspector to reveal the \textbf{Shape Popover}. It includes all of OmniGraffle's
    built-in shapes plus any custom shapes you've created.

    \ogfigure{0.800}{res/og7_inspectors24.png}
  \item Use the \textbf{Corner Radius} field to set how round the corners of the shape should
    be; enter 0 for perfectly pointy corners.
\end{itemize}

\begin{ogpro}
  \begin{itemize}
    \item In OmniGraffle Pro, with multiple shapes selected use the \textit{shape combination}
      buttons for crafting custom shapes. Available actions include \textbf{Combine Shapes},
      \textbf{Intersect Shapes}, \textbf{Subtract Shapes}, and \textbf{Uncombine Shapes}.

      \ogfigure{0.800}{res/og7_inspectors25.png}
  \end{itemize}
\end{ogpro}

\subsubsection{Using Adjustable Shapes}

Underneath the \textbf{Common Shapes} found in the Shape Popover, you'll see a row of five
\textbf{Adjustable} shapes that you can choose from:

\begin{itemize}
  \item Adjustable Arrow
  \item Adjustable Double Arrow
  \item Adjustable Arc
  \item Adjustable Wedge
  \item Adjustable Star
\end{itemize}

\ogfigure{0.800}{res/og7_insp_shapes_adjustable01.png}

As their name implies, adjustable shapes can be morphed to take on different appearances. The
thing that makes an Adjustable Shape \textit{adjustable} are the \textbf{Blue Adjustment
Handles}.

Depending on the shape, an Adjustable Shape has either one or two \textbf{Blue Adjustment
Handles}. The Adjustable Arrow and Double Adjustable Arrow have a single Adjustment Handle,
which you can use to resize the shape of the arrowhead.

\ogfigure{0.800}{res/og7_insp_shapes_adjustable02.png}

The \textbf{Adjustable Arc} and \textbf{Adjustable Wedge} shapes have two Adjustment Handles,
which you can grab and drag to change the degrees in the arc. If you hold down the
\textbf{Shift} key while adjusting the an Arc or Wedge, the shape increases or decreases in
size in 15-degree increments.

\ogfigure{0.800}{res/og7_insp_shapes_adjustable03.png}

The \textbf{Adjustable Arcs} and \textbf{Wedges} can be Filled with any color or pattern, and
arranged to create exploding pie charts.

\ogfigure{0.800}{res/og7_insp_shapes_adjustable04.png}

\textbf{Adjustable Stars} have two adjustment handles, which you can use to make the points more
or less pointier, and to decrease or increase the number of points on a star.

\ogfigure{0.800}{res/og7_insp_shapes_adjustable05.png}

\textbf{Adjustable Shapes} can really speed up your time to creation, so look for ways to use
them for diagrams, charts, wireframes, or website mockups.

\subsection{Defining Line Attributes and Behaviors with the Line Inspector}

Use this inspector to configure the properties of lines used to connect objects.

\ogfigure{0.800}{res/og7_inspectors26.png}

\begin{itemize}
  \item Use the pop-up menu at the top to set how the line travels from its source to its
    destination:

    \begin{itemize}
      \item \textbf{Straight} for a line that takes the shortest possible path between points.
      \item \textbf{Curved} for a line that turns smoothly to pass through each point.
      \item \textbf{Orthogonal} for a line that always travels horizontally or vertically.
      \item \textbf{Bézier} for a line with control points you can adjust with the selection tool.
    \end{itemize}
  \item In the middle of the Line inspector, you'll find two pop-up menus on the right and
    left---used to set the appearance of the tail and head of the line. Click the reverse button
    to swap the line's source and destination points.
  \item Use the fields below the tail and head menus to change the line ending sizes. If you've
    set the line type to \textbf{Orthogonal}, you can use the curve radius field in the middle
    to add some curvature to the line's corners.
  \item The \textbf{Line Hops} menu determines what the line should do when it crosses (or is
    crossed by) another line.

    To better understand how Line Hops work, it's important to understand how the things you draw
    or place on a layer stack in relation to each other.

    Everything you draw or place on the same layer of a canvas---including lines---shares the same
    \textit{stacking order}. The first thing drawn is at the bottom of that layer's object
    stack; all subsequent objects and lines drawn on that layer appear further up in the object
    stack.

    Consider the following:

    \ogfigure{0.800}{res/og7_linehops01.png}

    Use the \textbf{Line} tool to connect the objects, using the numbers on the lines shown here
    as the order in which to draw the lines:

    \ogfigure{0.800}{res/og7_linehops02.png}

    As you can see, Lines 3 and 5 cross each other. Following the principles of object stacking,
    Line 3 appears under Line 5. If you select Lines 3 and 5 and choose \textbf{Round hops over}
    in the \textbf{Line Hops} pop-up menu, Line 5 hops over Line 3 because it is lower in the
    stack.

    \ogfigure{0.800}{res/og7_linehops03.png}

    Conversely, if you choose \textbf{Round hops under} in the \textbf{Line Hops} pop-up menu,
    Line 3 loops underneath Line 5 because it is higher in the stack.

    \ogfigure{0.800}{res/og7_linehops04.png}

    If you would like a line to appear higher or lower in the object stack, select the line in
    question, choose the \textbf{Arrange} menu, and then choose one of the following menu
    options:

    \begin{itemize}
      \item Bring to Front (\textbf{Shift-Command-F})
      \item Bring Forward (\textbf{Option-Command-F})
      \item Send to Back (\textbf{Shift-Command-B})
      \item Send Backward (\textbf{Option-Command-B})
    \end{itemize}
\end{itemize}

\subsection{Placing Graphics Inside a Shape with the Image Inspector}

You can use this inspector to place and position an image inside a shape, or to work with an
image you have dragged directly to the canvas.

\ogfigure{0.800}{res/og7_inspectors27.png}

\begin{itemize}
  \item Click \textbf{Set Image} to choose and place an image in the selected shape. When you
    browse for an image file, clicking \textbf{Option} in the lower-left of the selection sheet
    provides you with options to \textbf{Include image in document} or \textbf{Link to image
    file}.

    \ogfigure{0.800}{res/og7_inspectors28.png}

    If you choose the \textbf{Link to image file} option, OmniGraffle places an alias of the image
    in the shape and relies on having access to that image file locally.

    \begin{ogtip}
      If you delete the linked image file from your Mac, or try to open the OmniGraffle file on
      another Mac, OmniGraffle displays an alert to let you know that it can't find the image. The
      alert displays the name of the image file and the linked location on your filesystem, so you
      can try to relocate or re-link the image.
    \end{ogtip}

    If an object's image comes from an alias, some extra commands become available in the Gear
    menu \ogicon{12.5pt}{res/og7_button_imageinsp_gear.png}:

    \begin{itemize}
      \item \textbf{Reveal in Finder} open a Finder window and takes you to where the image file
        is located.
      \item \textbf{Embed in Document} places a copy of the image file in the OmniGraffle file,
        allowing you to safely send the file to another OmniGraffle user.
      \item \textbf{Edit LinkBack Item} gives you the ability to edit a LinkBack item's resource.
    \end{itemize}
  \item Once an image is set, a \textbf{Remove} button becomes available so you can remove the
    image and try another.
\end{itemize}

Beneath the image well, you'll find the following controls for working with the image:

\begin{itemize}
  \item Use the slider to change the image's opacity:

    \ogfigure{0.800}{res/og7_inspectors29.png}
  \item The following row of buttons lets you choose how the image is placed within the shape.
    You can choose to use the image's natural size (left), to stretch the image to fit within
    the shape's dimensions (center), or to tile the image within the shape (right).

    \ogfigure{0.308}{res/og7_inspectors30.png}
  \item The positioning offsets are enabled if you are displaying the image at its natural
    aspect ratio, as set by the first of the three buttons to the left. You can enter a positive
    or negative percentage, or drag the image around in the image well to indicate the area you
    want.
  \item The zoom controls are available unless you are stretching the image to fit the shape.
    You can enter a zoom percentage, or just drag the slider, to make the image larger or
    smaller.

    \ogfigure{0.800}{res/og7_inspectors31.png}
  \item When you press the \textbf{Mask} button, a \textbf{Masking Shape} with an inner grid is
    placed on top of the image, and opaque corner brackets with lines in between are placed
    outside the masking area for resizing the image itself:

    \ogfigure{0.800}{res/og7_inspectors32.png}

    By default, the Masking Shape is a rectangle that's placed on top of the image. You can change
    the shape by choosing another from the Shape Popover in the \textbf{Shape} inspector (see
    Choosing a Shape with the Shape Inspector).

    \begin{itemize}
      \item To resize the Masking Shape, click-and-drag on the shape's bounding box handles.
      \item To resize the image, click-and-drag on the opaque image handles that are outside of
        the image itself.
      \item To reposition the image within the Masking Shape, you can either click and drag the
        image on the canvas, or use the image well to move the image around within the shape's
        area.
      \item When resizing the Masking Shape or the image itself, use the \textbf{Shift} key to
        proportionally resize the item, or use \textbf{Option-Shift} to proportionally resize the
        item based on its centerpoint.
      \item You can enter values in the \textbf{Offset} fields (horizontally and/or vertically),
        or to adjust the \textbf{Scale} values for the image's width and height.
    \end{itemize}

    When you've finished making adjustments to the image mask, click the \textbf{Mask} button
    again to hide the masking handles.
\end{itemize}

\begin{ogtip}
  While not exactly the same as cropping an image, using an image mask is a great way to get a
  crop effect by making only a selected portion of the image visible.
\end{ogtip}

At the bottom of the image inspector is a list of metadata associated with the image file,
including the following properties:

\begin{itemize}
  \item \textbf{Kind} --- The file type of the image.
  \item \textbf{Image Size} --- The horizontal and vertical dimensions of the original image, in
    pixels.
  \item \textbf{Visible Size} --- The horizontal and vertical dimensions of the image as
    displayed on the canvas, in pixels.
  \item \textbf{Resolution} --- The resolution of the original image, in dots per inch (DPI).
  \item \textbf{File Size} --- The file size of the image.
  \item \textbf{Color Space} --- The \href{https://en.wikipedia.org/wiki/Color_space}{color
    space} of the image.
\end{itemize}

\subsection{Choosing a Font, Size, and Color with the Font Inspector}

Use this inspector to modify the font properties of the selected objects, including style, size,
and color.

\ogfigure{0.800}{res/og7_inspectors33.png}

\begin{itemize}
  \item Click the \textbf{Font Family} menu to choose a font, or to open macOS's \textbf{Fonts
    Panel}.

    \ogfigure{0.800}{res/og7_inspectors34.png}

    \ogfigure{0.800}{res/og7_inspectors35.png}
  \item Enter a number to set the point size of the text.
  \item Choose a Typeface, if there is one available for the font family you've chosen.
  \item Click the color well to choose a font color.
\end{itemize}

At the bottom of the Font inspector, you will find controls for adjusting the Kerning (the space
between individual characters) and Leading (the space between lines of text).
\textbf{Auto-kern} is turned on by default, but you can uncheck that box if you so desire.

\subsection{Styling Text with the Text Position Inspector}

Use this inspector to control the positioning and appearance of the selected shape's text. If
you have OmniGraffle Pro, the Text Position inspector also includes \textit{offset
positioning} controls.

\ogfigure{0.800}{res/og7_inspectors36.png}

The controls in the Text Position inspector are as follows:

\begin{itemize}
  \item \textbf{Vertical Text Alignment} --- The first three buttons are used for setting the
    vertical alignment of the text; choose to align to the top, middle, or bottom of the shape.
  \item \textbf{Horizontal Text Alignment} --- With the horizontal alignment buttons, you can
    align left, center, align right, or justify the text.
  \item \textbf{Text Positioning} --- There are five buttons for controlling how the text
    relates to the shape; you can choose:

    \begin{itemize}
      \item To fill the shape with the text
      \item To let the text overflow outside of the shape
      \item To crop off the text that doesn't fit inside the shape
      \item To resize the shape to make the text fit
      \item To make the text follow the shape's path
    \end{itemize}

    If you choose to have the shape resize, the object will only be manually resizable in the
    horizontal direction; it will resize vertically according to its text.
  \item Select the \textbf{Wrap to Shape} checkbox to limit the text's width to the width of the
    shape.
\end{itemize}

\subsubsection{Adding Text to the Path of a Shape}

To add text as the border of a shape, double-click the shape to enter the text, then choose
\textbf{Text Follows Shape's Path} (the farthest right Text Positioning button).

\ogfigure{0.800}{res/og7_text_on_a_shape.png}

\begin{ogtip}
  To add text to the path of a line instead, see Adding Text to the Path of a Line.
\end{ogtip}

\subsubsection{Offsetting Text in the Text Position Inspector
(\texorpdfstring{\textbf{PRO}}{PRO})}

\begin{ogpro}
  OmniGraffle Pro includes these additional text offset and positioning controls:

  \begin{itemize}
    \item By default, the \textbf{Use default offsets} option is turned on. If you uncheck this
      option, you can adjust the text's position using the four input fields:

      \begin{itemize}
        \item \textbf{Text Offset from Left} --- This adjusts the text's position from the left
          edge of the object.
        \item \textbf{Text Offset from Top} --- This adjusts the text's position from the top edge
          of the object.
        \item \textbf{Width} --- This adjusts the text box's width.
        \item \textbf{Height} --- This adjusts the text box's height.
      \end{itemize}
    \item Use the text's rotation value to determine how the text appears when you rotate the
      object:

      \begin{itemize}
        \item \textbf{Relative} is turned on by default, which means that the text within an
          object will maintain its relative position when an object is rotated.
        \item \textbf{Absolute} permits the text to be rotated independent of its bounding object.
      \end{itemize}

      You can either enter a rotation value in the field or drag the circular control to rotate
      the text.
  \end{itemize}
\end{ogpro}

\subsection{Changing Object Position and Size with the Geometry Inspector}

This inspector deals with an object's position and orientation.

\ogfigure{0.800}{res/og7_inspectors37.png}

\begin{itemize}
  \item At the top of the inspector, you'll find a pop-up menu with the following options:

    \begin{itemize}
      \item \textbf{No Scaling} --- by default, the objects you draw on the canvas are not scaled
        to any size or proportion.
      \item \textbf{Maintain Aspect Ratio} --- the objects you draw or resize are scaled
        proportionally. In addition, there are three additional aspect-related options:

        \begin{itemize}
          \item \textbf{Scale Stroke} --- scales the object's stroke proportionally.
          \item \textbf{Scale Font} --- proportionally scales any text inside an object, while also
            scaling the text object.
          \item \textbf{Scale Stroke and Font} --- scales the objects stroke and text
            proportionally.
        \end{itemize}
    \end{itemize}
  \item Underneath the scaling pop-up menu are two fields which display the X and Y (horizontal
    and vertical) coordinates for the selected objects; the value represents the distance
    between the canvas's origin and the upper-left corner of the object:

    \ogfigure{0.800}{res/og7_insp_geometry01.png}

    If you click on the blue arrow buttons next to the fields, you can measure the distance from
    the canvas origin to the top- or left-center of the object\ldots{}

    \ogfigure{0.800}{res/og7_insp_geometry02.png}

    \ldots{}or to the bottom-right corner of the selected object:

    \ogfigure{0.800}{res/og7_insp_geometry03.png}

    If you have multiple objects selected, two hyphens (\textbf{--}) appear in the fields that
    have different values. Enter a new X or Y value to move an object on the canvas.
  \item The next two fields contain the Width and Height values for the selected objects. If you
    have multiple objects selected, two hyphens (\textbf{--}) appear in the fields that have
    different values.

    To change the dimensions of an object, enter a new Width or Height in the fields. Entering a
    number value and pressing Tab or Return changes the Width or Height based on the
    \textbf{Ruler Units} set in the Units inspector. However, you can enter a value for
    \textit{any} unit of measurement, even percentages. OmniGraffle takes the value you have
    entered, and converts that into the unit of measurement you have specified in the Units
    inspector. For example, if \textbf{Ruler Units} is set to pixels (px) and you enter
    \textbf{2 cm} into the Width field, OmniGraffle changes the width of the object to
    \textbf{56.6929 px}.
  \item To rotate the object, enter a degree value in the rotation field, or drag the circular
    control.
  \item Click the flip buttons to reverse the object horizontally or vertically.
\end{itemize}

If you have selected a line label, two additional controls become active:

\ogfigure{0.800}{res/og7_inspectors38.png}

\begin{itemize}
  \item \textbf{Label Horizontal Position} --- Use the slider at the top of the inspector to
    adjust the label's position along the line, moving it to the left or right.
  \item \textbf{Label Orientation Menu} --- Use the pop-up menu to set how the label's text
    appears in relation to the line. Options include: Horizontal, Vertical, Parallel,
    Perpendicular, Independent, and Follows Path.
  \item \textbf{Label Vertical Position} --- Use the slider at the bottom of the inspector to
    position the label above or below the line, or somewhere in between.
\end{itemize}

\subsubsection{Adding Text to the Path of a Line}

To add text that follows the path of a line, select the Text tool and click the line to add a
line label. Enter the text, then choose \textbf{Follows Path} from the Label Orientation menu.

\ogfigure{0.800}{res/og7_text_on_a_line.png}

\begin{ogtip}
  To add text to the border of a shape instead, see Adding Text to the Path of a Shape.
\end{ogtip}

\subsection{Positioning Objects Neatly with the Alignment Inspector}

Use this inspector to uniformly align objects, either in relation to themselves, or to the
canvas. At the top of the inspector, choose between \textbf{Align to Shape} or \textbf{Align
to Canvas}.

\subsubsection{Align to Shape}

If you choose \textbf{Align to Shape}, the align and spread buttons consider the shape's
position in relation to other selected shapes. This makes it so you can precisely align and
set horizontal and vertical spacing between objects.

\ogfigure{0.800}{res/og7_inspectors39.png}

How other objects align is based on the position of the first object selected. For example, if
you have three objects on the canvas and you want the objects to align based on the position
of the center object, \textbf{Command-click} the center object first, and then
\textbf{Command-click} to select the other two objects:

\ogfigure{0.800}{res/og7_insp_aligntoshape01.png}

There are eight buttons for controlling how the selected objects will align to each other, as
follows:

\begin{itemize}
  \item \ogiconw{30.0pt}{res/og7_button_aligntop.png} aligns objects along their top edges
  \item \ogiconw{30.0pt}{res/og7_button_alignhcenter.png} aligns objects based on their
    horizontal center
  \item \ogiconw{30.0pt}{res/og7_button_alignbottom.png} aligns objects based on their bottom
    edges
  \item \ogiconw{30.0pt}{res/og7_button_alignleft.png} aligns objects based on their left edges
  \item \ogiconw{30.0pt}{res/og7_button_alignvcenter.png} aligns objects based on their vertical
    center
  \item \ogiconw{30.0pt}{res/og7_button_alignright.png} aligns objects based on their right
    edges
\end{itemize}

Off to the right are two additional positioning controls:

\begin{itemize}
  \item \ogiconw{30.0pt}{res/og7_button_spreadv.png} spreads the selected objects evenly based
    on the positions of the top and bottom objects
  \item \ogiconw{30.0pt}{res/og7_button_spreadh.png} spreads the selected objects evenly based
    on the positions of the left and right objects
\end{itemize}

\begin{ognote}
  Use Canvas Preferences to switch the default behavior of the Align to Shape controls between
  aligning to the \textit{first object} (as described above) and aligning objects to the
  \textit{selection} instead.

  When aligning to the selection, objects are positioned based on the location of the object
  furthest in the direction of the alignment choice, as follows:

  \begin{itemize}
    \item \ogiconw{30.0pt}{res/og7_button_aligntop.png} aligns objects along the top edge of the
      uppermost object
    \item \ogiconw{30.0pt}{res/og7_button_alignhcenter.png} aligns objects to the horizontal
      center of the uppermost object
    \item \ogiconw{30.0pt}{res/og7_button_alignbottom.png} aligns objects along the bottom edge
      of the lowermost object
    \item \ogiconw{30.0pt}{res/og7_button_alignleft.png} aligns objects along the left edge of
      the leftmost object
    \item \ogiconw{30.0pt}{res/og7_button_alignvcenter.png} aligns objects along the horizontal
      center of the leftmost object
    \item \ogiconw{30.0pt}{res/og7_button_alignright.png} aligns objects along the right edge of
      the rightmost object
  \end{itemize}
\end{ognote}

\subsubsection{Align to Canvas}

If you choose \textbf{Align to Canvas}, the align and spread buttons consider the canvas edges
as ``objects'' when arranging the selected objects. This makes it so that you can align
objects to a canvas edge, or space objects evenly between the canvas edges.

\ogfigure{0.800}{res/og7_inspectors40.png}

\begin{itemize}
  \item \ogiconw{30.0pt}{res/og7_button_canvastop.png} aligns objects to the top edge of the
    canvas
  \item \ogiconw{30.0pt}{res/og7_button_canvashcenter.png} aligns objects to the horizontal
    center of the canvas
  \item \ogiconw{30.0pt}{res/og7_button_canvasbottom.png} aligns objects to the bottom edge of
    the canvas
  \item \ogiconw{30.0pt}{res/og7_button_canvasleft.png} aligns objects to the left edge of the
    canvas
  \item \ogiconw{30.0pt}{res/og7_button_canvasvcenter.png} aligns objects to the vertical center
    of the canvas
  \item \ogiconw{30.0pt}{res/og7_button_canvasright.png} aligns objects to the right edge of the
    canvas
\end{itemize}

Off to the right are two additional positioning controls:

\begin{itemize}
  \item \ogiconw{30.0pt}{res/og7_button_canvasspreadv.png} spreads the selected objects evenly
    between the top and bottom edges of the canvas
  \item \ogiconw{30.0pt}{res/og7_button_canvasspreadh.png} spreads the selected objects evenly
    between the left and right edges of the canvas
\end{itemize}

\subsubsection{Aligning to the Grid}

If you are working on a grid layout or creating pixel art, you may opt to use either of the two
buttons on the right side of the Alignment inspector:

\begin{itemize}
  \item \ogiconw{30.0pt}{res/og7_button_aligngridcenter.png} aligns the selected objects to the
    grid's center.
  \item \ogiconw{30.0pt}{res/og7_button_aligngridedge.png} aligns the selected objects to the
    grid's edge.
\end{itemize}

\subsubsection{Controlling How Objects Space Apart}

At the bottom of the Alignment inspector are two controls for setting the spacing between
objects:

\begin{itemize}
  \item \ogiconw{30.0pt}{res/og7_button_alignspaceh.png} sets the spacing between objects
    horizontally
  \item \ogiconw{30.0pt}{res/og7_button_alignspacev.png} sets the spacing between objects
    vertically
\end{itemize}

Enter a value in the field and then click either the horizontal or vertical spacing button.
These buttons are only useful if you have at least two objects selected.

\subsection{Applying Object Styles with the Style Tray}

Beneath the Inspectors in the bottom bar you'll find the \textbf{Style Tray}, a row of Style
Swatches (also known as \textit{chits}) that provide a summary of the styles applied to the
currently selected object.

\ogfigure{0.746}{res/og7_diagramming_styletray.png}

From left to right, the Swatches represent:

\begin{itemize}
  \item Summary (All Styles)
  \item Fill
  \item Stroke
  \item Image
  \item Shadow
  \item Shape
  \item Font
  \item Text Position
\end{itemize}

Click and drag a Style Swatch from the tray onto an object to apply that style (or set of
styles) to it.

\begin{ogtip}
  You can also drag styles from the Style Tray straight onto items in the Sidebar's Outline or
  Layers tabs.
\end{ogtip}

\section{The Properties Inspectors}

Use the \textbf{Properties inspectors} to manage how objects connect to one another, and in
OmniGraffle Pro, set actions and add notes and key-value pair relationships to objects.

\subsection{Changing the Way Objects Connect with the Connections Inspector}

Use this inspector to control how lines and objects connect to one another.

\subsubsection{Connections To and From Lines}

Use the \textbf{Lines} tab of the Connections inspector to define the relationship of the lines
you draw with the \textbf{Line} tool.

\ogfigure{0.800}{res/og7_inspectors41.png}

\begin{itemize}
  \item \textbf{Allow connections from lines} permits lines to connect to each other.
  \item \textbf{Allow connections to other objects} defines whether the selected line(s) connect
    to other objects.
  \item \textbf{Allow shapes to become labels} lets you set a shape as the background to a line
    label.
  \item \textbf{Route from object centers} defines whether connecting lines between objects
    route to/from the object's center point or based on \textbf{Magnet} positions.
\end{itemize}

For example, if you are using the Line tool to draw in OmniGraffle, you may want to turn off
\textbf{Allow connections from lines}. If that option is turned on, the lines you draw will
try to connect to one another, which can be problematic when freehand sketching with lines.

\subsubsection{Connections with Other Objects}

Use the \textbf{Other Objects} tab to define how shapes accept connections from other objects on
the canvas.

\ogfigure{0.800}{res/og7_inspectors42.png}

\begin{itemize}
  \item Turn off \textbf{Allow connections from lines} to make it impossible for connection
    lines to use the selected objects as sources or destinations.
  \item If you have a group or a table selected, you can choose whether connection lines should
    be allowed to connect to objects in the group or only to the group as a whole.
  \item The \textbf{Magnet Positions} pop-up menu contains a bunch of magnet presets that you
    can choose and assign to the selected shape. Magnets are points on a shape that attract
    connection lines. You can choose to have \textbf{No magnets} (the default), magnets placed
    according to cardinal directions (North, South, East, and West), magnets on each vertex
    (corner), or a certain number of magnets on each side of the shape.

    If you select some number of magnets per side, then hold \textbf{Shift} while opening the
    pop-up menu and selecting another number, the two numbers are added together; you can get up
    to 10 magnets per side in this way.

    \begin{ogtip}
      Of course, you can always use the Magnet tool to add magnets to a shape. See Adding and
      Adjusting Connections with the Magnet Tool for more information about the Magnet tool.
    \end{ogtip}
  \item If you have a line selected, the Lines checkboxes become available. Deselect the
    \textbf{Allow connections to other objects} checkbox to make it impossible for the line to
    have an object as its source or destination. Deselect the \textbf{Allow shapes to become
    labels} checkbox to make it impossible to drag a shape onto the line and make it a line
    label; existing labels stay attached.
  \item Use the \textbf{Rank} pop-up menu to set objects' \textit{preferred rank}, which affects
    how the Diagram Layout feature positions the objects relative to any connected objects.

    \begin{itemize}
      \item \textbf{Default Object Ranking} lets OmniGraffle decide the rank based on the objects'
        original position in the hierarchy.
      \item \textbf{Minimum Rank} places the objects as peers of the top level of the hierarchy.
      \item \textbf{Maximum Rank} places the objects as peers of the bottom level of the
        hierarchy.
      \item \textbf{Same Rank} ensures that the objects end up on the same level (only available
        when multiple objects are selected).
    \end{itemize}

    These assignments don't change the relationships between objects (though they may appear to
    shift in the outline), so you can always select all of your objects and choose
    \textbf{Default Object Ranking} to return them to normal.
\end{itemize}

\begin{ogtip}
  Setting object rank is particularly useful for fine-tuning a diagram where certain relationships
  may not be possible to depict using the outline alone.

  For example, show a marriage between members of two houses by using the Line tool to connect
  the two objects that represent them, select both objects, and choose \textbf{Same Rank} from
  the pop-up menu.
\end{ogtip}

\subsection{Making Objects Interactive with the Action Inspector
(\texorpdfstring{\textbf{PRO}}{PRO})}

\begin{ogpro}
  The \textbf{Action inspector} determines what should happen when someone clicks the selected
  object with the Action Browse Tool \ogiconw{19.4pt}{res/og7_tools_actionbrowse.png} in the
  toolbar. The default action is \textbf{Does Nothing}. You can click the object all you want,
  and nothing happens.

  \begin{itemize}
    \item \textbf{Open a URL} --- Specify a URL or filepath to open. Links (URLs) to web sites
      are opened in your default web browser, while files are opened in the appropriate app.

      \ogfigure{0.800}{res/og7_inspectors43.png}
    \item \textbf{Open a File} --- Choose a file to open; when the Action is clicked, the file
      will open in its default app.

      \ogfigure{0.800}{res/og7_inspectors44.png}

      Note that file paths are absolute: they start from the root \texttt{Users} folder of your
      hard drive. Click \textbf{Open} button to try opening the file or URL that you've
      specified. When the object is clicked, the file or URL is opened in the appropriate app.
    \item \textbf{Run AppleScript} --- Runs an AppleScript that you enter into the space
      provided.

      \ogfigure{0.800}{res/og7_inspectors45.png}

      If you choose \textbf{Run AppleScript}, you get a text field for entering an AppleScript.
      The script that you enter is run such that \texttt{self} refers to the clicked object.
      Click \textbf{Check Syntax} to make sure that the AppleScript is correct, and then click
      \textbf{Run Script} to try it out. In Presentation Mode (\textbf{Option-Command-P}), the
      script runs when an actionable object is clicked.
    \item \textbf{Run Automation Action} --- Runs an Omni Automation Plug-In, which you specify
      by selecting the script file from the \textbf{Automation} menu. After choosing the
      Plug-In, you can assign an \textbf{Action} to the script.

      \ogfigure{0.800}{res/og7_insp_action_automation.png}

      To learn more about writing Omni Automation Plug-Ins with JavaScript, visit the
      \href{https://www.omni-automation.com}{Omni Automation website}.
    \item \textbf{Jump to} --- Jump to another location in the current OmniGraffle file. You can
      specify a specific canvas name, and choose the location type from the pop-up menu that
      appears above the preview.

      \ogfigure{0.800}{res/og7_inspectors46.png}

      Locations you can choose from include:

      \begin{itemize}
        \item Highlight an Object
        \item Center on a Point
        \item Zoom to Display a Rectangle
        \item Switch to the Previous Canvas
        \item Switch to the Next Canvas
        \item Switch to a Specific Canvas
      \end{itemize}

      Some of these options offer a canvas preview, in which you can click or drag to indicate
      which object to highlight, which point to center on, or where to zoom. If you select
      \textbf{Switch to a Specific Canvas}, use the \textbf{Canvas} pop-up menu to choose from
      one of the canvases in your project.
    \item \textbf{Show or Hide Layers} --- You can indicate whether to Show, Hide, or Toggle the
      visibility of any layer of the current canvas.

      \ogfigure{0.800}{res/og7_inspectors47.png}
  \end{itemize}
\end{ogpro}

\subsection{Assigning Metadata to Objects With the Object Data Inspector
(\texorpdfstring{\textbf{PRO}}{PRO})}

\begin{ogpro}
  The Object Data inspector contains a field for entering text to associate with the selected
  object, and a table of custom data.

  \ogfigure{0.800}{res/og7_inspectors48.png}

  Object Data can be formatted as Rich Text, which means you can use all of the different font
  styles and colors that you can use elsewhere in OmniGraffle. When you place the mouse
  pointer over an object on the canvas with data assigned, the object's data appears in a help
  tag floating over the object.

  You can use the custom data table to keep your own information about the object. Custom data
  is stored as key/value pairs: the \textbf{Data Key} is like a label for what type of
  information you are storing, and the \textbf{Data Value} is the information itself.

  For example, imagine you have a diagram of a computer network, and you want to assign a model
  number to each component. Click the placeholder row or the plus button to create a new
  key/value pair. In the Data Key column, you would type Part Number, and in the Data Value
  field, you would type, say, A1181. This data doesn't have any effect on the way OmniGraffle
  works; it's just a way for you to store pertinent data about objects in your diagram. To
  delete a row of data, click the `x' button on the right side of the row.

  The information you add to the Object Data inspector is also exposed to macOS's Spotlight
  search feature, in case you need to find which of your OmniGraffle diagrams contains
  specific information.
\end{ogpro}

\section{The Canvas Inspectors}

Use the Canvas inspectors to manage the appearance and properties of the current or selected
canvases.

\subsection{Define the Size and Autosizing Properties of Your Canvas with the Canvas Size
Inspector}

Use the Canvas Size inspector to specify your type of canvas; choose between:

\begin{itemize}
  \item Fixed
  \item Flexible
  \item Infinite
\end{itemize}

\subsubsection{Fixed}

The canvas is a size of your specification, based on the values set in the \textbf{Width} and
\textbf{Height} fields. Choose a Fixed Canvas Mode when you know the dimensions of the thing
you are creating; for example, an app icon.

\ogfigure{0.800}{res/og7_insp_canvassize_fixed.png}

\begin{itemize}
  \item \textbf{Width}: Enter a different value or click the tiny arrows to decrease or increase
    the width of the canvas.
  \item \textbf{Height}: Enter a different value or click the tiny arrows to decrease or
    increase the height of the canvas.
\end{itemize}

\subsubsection{Flexible}

A canvas that can be expanded in any direction: up, down, left, right, or any combination of the
four. Choose a Flexible Canvas Mode when want the canvas to expand or collapse to the size of
objects on the canvas.

\ogfigure{0.800}{res/og7_insp_canvassize_flexible.png}

\begin{itemize}
  \item \textbf{Width}: Enter a different value or click the tiny arrows to decrease or increase
    the width of the canvas. Changing the Width after the Autosizing Switches have been set
    turns off the left and right switches.
  \item \textbf{Height}: Enter a different value or click the tiny arrows to decrease or
    increase the height of the canvas. Changing the Height after the Autosizing Switches have
    been set turns off the top and bottom switches.
\end{itemize}

Use the \textbf{Autosizing Switches} to enable canvas expansion in that direction.

\ogfigure{0.800}{res/og7_canvas_autosizeswitches.png}

\begin{itemize}
  \item If your design has a specific width, you can set the \textbf{Canvas Width}, and then
    choose which direction the canvas autosizes vertically (Up, Down, or both).
  \item Similarly, if your design has a specific height, you can set the \textbf{Canvas Height},
    and then choose which direction the canvas autosizes horizontally (left, right, or both).
\end{itemize}

\subsubsection{Infinite}

A canvas that infinitely extends in all directions. Choose an Infinite Canvas Mode when you have
a document that needs a lot of space.

\ogfigure{0.800}{res/og7_insp_canvassize_infinite.png}

\begin{ognote}
  The \textbf{Width} and \textbf{Height} of an Infinite canvas are not adjustable.
\end{ognote}

Additional options you'll find in the Canvas Size inspector include:

\begin{itemize}
  \item With \textbf{Show page breaks} selected, lines appear on the canvas to indicate how your
    document will be printed.
  \item With \textbf{Size uses printer pages} selected, the Canvas Size is determined by the
    \textbf{Paper Size} you have selected in the \textbf{Print Setup} screen (\textbf{File}
    $\blacktriangleright$ \textbf{Print Setup}, or \textbf{Shift-Command-P}).
  \item The \textbf{Canvas Scroll Margin} field lets you adjust how much area is added as a
    scrolling margin around the canvas. When set to \textbf{100\%}, your entire canvas can
    scroll out of view.
  \item The \textbf{Orientation} of pages can be \textbf{Portrait} (vertical),
    \textbf{Landscape} (horizontal), or can be taken from the settings in \textbf{Page Setup},
    which is the default setting.
  \item With \textbf{Print canvas on one page} selected, the canvas gets scaled up or down when
    printed to fit on a single piece of paper, regardless of its size on the screen.
\end{itemize}

\begin{ogpro}
  \begin{itemize}
    \item If you have OmniGraffle Pro, you can also specify how things look onscreen when
      viewing your document at \textbf{100\% Zoom}:

      \begin{itemize}
        \item \textbf{One screen pixel} matches the pixel density of your display. For example, at
          100\% zoom, objects viewed at \textbf{One screen pixel} only appear smaller onscreen;
          however, they retain their actual size (which you can verify by looking at the object's
          dimensions in the \textbf{Geometry} inspector).
        \item \textbf{One Apple point} (the default) relates to resolution-independent distance on
          the display. For example, on a standard-resolution display, 1 point equals 1 pixel; on a
          Retina display, 1 point equals 2 pixels. At 100\% zoom, objects viewed at \textbf{One
          Apple point} only appear smaller onscreen; however, they retain their actual size (which
            you can verify by looking at the object's dimensions in the \textbf{Geometry}
          inspector).
        \item \textbf{One PostScript point} displays the contents of your document at real-world
          size. At 100\% zoom, you can hold a ruler up to your display and the size of the object
          you are measuring onscreen will match the dimensions you see in the \textbf{Geometry}
          inspector.
      \end{itemize}

      The \textbf{100\% Zoom} options make it possible for you to work on large-scale projects
      (such as a poster for an event, or even a highway billboard) at a smaller scale onscreen,
      and provide you with the flexibility to view your work in real-world dimensions when
      needed.
  \end{itemize}
\end{ogpro}

\subsection{Specify a Background Fill or Pattern in the Canvas Fill Inspector}

Select a canvas by clicking its preview in the sidebar to edit its background fill properties.
As with the Fill inspector (see Coloring Shapes with the Fill Inspector, earlier in this
chapter), you can choose from nine fill styles and edit various parameters of each.

\ogfigure{0.800}{res/og7_inspectors51.png}

\subsection{Use an Image as a Backdrop in the Background Image Inspector}

You can also add an image to a canvas. Just click the canvas in the sidebar and then use the
Background Image inspector to place an image similarly to the Image inspector (see Placing
Graphics Inside a Shape with the Image Inspector, earlier in this chapter).

\ogfigure{0.800}{res/og7_inspectors52.png}

Use the controls above the Mask button to Manually Size, Stretch, or Tile the image you've
placed. Use the grid of four fields to the right to position the image from left, from top,
sized horizontally, or sized vertically. Use the slider beneath the image well to change the
image's opacity.

\subsection{Defining the Units of Measurement in the Units Inspector}

Use the \textbf{Units inspector} to determine the unit of measure to use for the canvas and its
rulers, as well as to set the scale and ruler origin points.

\ogfigure{0.800}{res/og7_inspectors53.png}

\subsubsection{Ruler Units}

The \textbf{Ruler Units} pop-up menu is where you set the type of measurement units you'd like
to use for the current canvas.

\ogfigure{0.781}{res/og7_inspectors54.png}

\begin{ognote}
  You can also change the Ruler Units by Control- or Right-clicking on the Canvas or one of the
  rulers. For more information, see Changing Ruler Units to Suit Your Design Needs earlier in
  this guide.
\end{ognote}

The actual size of your diagram does not change when you change the Ruler Units; it is merely
measured differently. The ruler and the inspectors display measurements in whichever unit you
select.

As you may have noticed, \textit{Points} (pt), \textit{Centimeters} (cm), and \textit{Inches,
fractional} (in) each have an asterisk (\textbf{*}) next to them in the Ruler Units pop-up
menu. With one of these Ruler Units selected, objects on the canvas will print to size as long
as the Scale is 1:1 (i.e., \textbf{1 pt = 1 pt}, \textbf{1 cm = 1 cm}, or \textbf{1 in = 1
in}). For example, if you have a 5 cm square on the canvas and you print that canvas (by
choosing \textbf{File} $\blacktriangleright$ \textbf{Print}), you could measure that square
with a ruler to reveal it is the exact size and scale as noted in the Units inspector.

\subsubsection{Scale}

The \textbf{Scale} pop-up menu can be used to interpret simple expressions of scale. This allows
you, for example, to use OmniGraffle for things like creating to-scale drawings of your house
and garden area, or for scaling icons and graphics needed for an app or website.

By default, the scale is based on what you have selected in the \textbf{Ruler Units} pop-up. For
example, if you set the Ruler Units to \textbf{feet (ft)}, the \textbf{Scale} pop-up reads
\textbf{12 in = 1 ft}.

\ogfigure{0.800}{res/og7_inspectors55.png}

If you choose \textbf{Custom} from the \textbf{Scale} pop-up, you can change the scale to suit
the project you're currently working on. For example, if you enter \textbf{1 cm = 1 m}, 1 cm
on the ruler now becomes 1 meter, objects on the canvas that were 2 cm wide are now 2 meters
wide, and so on. The Units setting changes to match the second value in the equation.

You can also enter a ratio, such as \textbf{1:2} or \textbf{100:93000000}, as a \textbf{Custom}
scaling option. The Scale ratio you set respects the Ruler Units selected in the pop-up menu
above. For example, if the Ruler Units are set to \textbf{pixels (px)}, and the scale uses a
ratio of \textbf{1:2}, objects on the Canvas are, effectively, twice as big.

\subsubsection{Enforce Canvas Scale}

If you need to draw something to a different scale (for example, if you are doing landscape
planning), clearly you don't want to print an exact-scale replica of your garden. However, if
you first select \textbf{Enforce canvas scale} in the Units inspector, you can change the
scale to something more sane; such as \textbf{1 in = 12 in}, where 1 inch on the canvas is the
equivalent of 12 inches (1 foot) in the real world. You'll notice that the objects and rulers
change their scale, but the object retains its dimensions in the \textbf{Geometry} inspector.

Optionally, you can turn on \textbf{Enforce canvas scale} in the Units inspector. With
\textbf{Enforce canvas scale} enabled, nothing happens to the objects on the Canvas when the
Scale is at a 1:1 ratio (e.g., \textbf{1 pt = 1 px} or \textbf{1 in = 1 in}). If you change
the Scale ratio to 1:3, for example, the objects appear to be smaller as the canvas scales to
the new setting. However, if you select an object and examine its properties in the
\textbf{Geometry} inspector, you'll notice that the object is still the same size---it just
looks smaller because the Canvas Scale's is no longer 1:1.

The one caveat to enabling \textbf{Enforce canvas scale} is when you are adding stencils to your
canvas. Stencils, such as the \textbf{Basic} $\blacktriangleright$ \textbf{Shapes} objects,
already have a predetermined size:

\ogfigure{0.800}{res/og7_inspectors56.png}

With a 1:1 Scale---regardless of whether \textbf{Enforce canvas scale} is enabled---stencils you
place on the canvas are sized normally. If you turn on \textbf{Enforce canvas scale} and then
change the ratio to something like 1:4, the stencil appears to be smaller, but it is sized to
fit the Canvas Scale. However, if you have the same 1:4 canvas scale ratio but turn off
\textbf{Enforce canvas scale}, the stencil not only appears to be four times larger, it
\textit{is} four times larger when you examine its properties in the \textbf{Geometry}
inspector.

\subsubsection{Canvas Origin Points}

By default, the very upper-left corner of a canvas is its origin (that is, the point where the
rulers' measurements start from, where the coordinates are 0,0). To change the origin, enter
values in the two Origin fields. (You can also drag the origin from the corner where the
rulers meet.) The coordinates in the Geometry inspector are based on this origin point.

\begin{ognote}
  In the case of an Infinite Canvas, a base \textit{point} on the canvas is used as the origin
  (0,0). The numbers on the rulers can extend in all directions, from negative to positive. See
  What Happens to the Canvas Origin with an Infinite Canvas, earlier in this guide.
\end{ognote}

\subsection{Setting Up a Grid with the Grid Inspector}

Use this inspector to set up a grid on the canvas, so you can keep objects lined up nicely.

\ogfigure{0.800}{res/og7_inspectors57.png}

The two fields at the top of the inspector are where you can set the values for the Major and
Minor Grid spacing, respectively:

\begin{itemize}
  \item The \textbf{Major Grid Spacing} field determines how large each square of the main grid
    should be. Click the color well next to it to choose a color for the major grid.
  \item The \textbf{Minor Grid Steps} field determines how many minor grid squares should fit
    across one major grid square. For example, if your major grid squares are 100 pixels across,
    and you have 10 minor grid steps, your minor grid squares are 10 pixels across. Click the
    color well next to it to choose a color for the minor grid.
\end{itemize}

The checkboxes beneath the Major and Minor Grid spacing fields provide additional control over
how the grid is used and its appearance:

\begin{itemize}
  \item While the \textbf{Snap to grid} checkbox is selected, all objects fit themselves to the
    grid as you create or move them. (If you want to snap objects already on the canvas, use the
    Align To Grid button below.)
  \item Select \textbf{Show grid lines} to show the Minor Grid lines on the canvas.

    \begin{itemize}
      \item Select \textbf{Show major} to make the heavier Major Grid lines visible as well.
      \item Select \textbf{Grid in front} to make the grid visible in front of objects on the
        canvas.
      \item Select \textbf{Print grid} to include the grid when printing.
    \end{itemize}
\end{itemize}

\begin{ogtip}
  Choose \textbf{Arrange} $\blacktriangleright$ \textbf{Grid} $\blacktriangleright$ \textbf{Align
  Objects to Grid} (\textbf{Option-Command-[}) to make all of the selected objects line up to
    the grid right away.
  \end{ogtip}

  \subsubsection{Nudging Objects on the Canvas}

  While working in OmniGraffle, there are times when you just need to give an object a little
  ``nudge'' to get it into position. To nudge an object, select it first and then press one of
  the arrow keys on the keyboard to move it up, down, left, or right.

  The \textit{nudge distance} is determined by the distance set in the Grid inspector's
  \textbf{Minor Grid Steps} field, and whether \textbf{Snap to grid} is on or off:

  \begin{itemize}
    \item \textbf{When Snap to Grid is Off}, pressing an arrow key nudges the selected object 1 pt
      in the direction of the arrow key pressed, regardless of the units you are using.
    \item \textbf{When Snap to Grid is On}, pressing an arrow key nudges the selected object by
      one \textbf{Minor Grid Step} in the direction of the arrow key pressed.
  \end{itemize}

  \begin{ogtip}
    Holding \textbf{Shift} while nudging inverts the nudge behavior. For example, if \textbf{Snap to
    grid} is \textit{off} and you hold \textbf{Shift} while nudging an object, the object is
    nudged by one \textbf{Minor Grid Step}.

    Similarly, if \textbf{Snap to grid} is \textit{on} and you hold the \textbf{Shift} key while
    nudging an object, the object is nudged by 1 pt in the direction of the arrow key pressed.
  \end{ogtip}

  \subsection{Automatically Arranging and Styling Objects with the Diagram Layout and Style
  Inspector}

  Use this inspector to automatically lay out shapes based on the logical relationships
  established by the connection lines between them, and to apply a diagram style to connected
  objects on the canvas.

  Use the \textbf{Include} pop-up menu to determine whether auto layout and diagram style
  properties apply to all objects on the canvas, or only objects that are connected to one
  another.

  Use the \textbf{Layout Type} pop-up menu to select from one of four layout types: Hierarchical
  (the default), Force-directed, Circular, and Radial. The various controls within the Diagram
  Layout section of the inspector change depending on which layout type you choose.

  \begin{description}[style=nextline]
    \item[{Hierarchical}] The hierarchical layout creates layers of equally-ranked objects,
      extending in one direction.

      \ogfigure{0.800}{res/og7_inspectors58.png}
    \item[{Force-directed}] The force-directed layout grows in semi-random directions from the
      center, rather than in one particular direction from the edge.

      \ogfigure{0.800}{res/og7_inspectors59.png}

      The top slider determines the length of the connecting lines between objects, and the bottom
      slider determines how far apart the objects appear in the diagram.
    \item[{Circular}] The circular layout tries to arrange sibling shapes in a circle around their
      parent.

      \ogfigure{0.800}{res/og7_inspectors60.png}
    \item[{Radial}] The radial layout tries to arrange sibling shapes in arcs around their parent.

      \ogfigure{0.800}{res/og7_inspectors61.png}
  \end{description}

  When you use a hierarchical layout:

  \begin{itemize}
    \item The Direction buttons change where the top level objects start, and which way to layer
      the lower level objects from there.
    \item The Rank Separation field controls how far away each level of objects should be from the
      next.
    \item The Object Separation field controls how far away each object should be from other
      objects on the same level.
  \end{itemize}

  When you use other layout methods:

  \begin{itemize}
    \item Connection lines can stretch and compress, but you can adjust their average length by
      dragging the Line Length slider.
    \item The Shape Repulsion slider determines how strongly the shapes try to avoid coming near
      one another. If the line length and shape repulsion are small enough, shapes can be made to
      overlap.
  \end{itemize}

  Finally, you can turn on \textbf{Auto layout} to make OmniGraffle distribute the objects on the
  canvas whenever the connections between them change.

  \ogfigure{0.800}{res/og7_inspectors66.png}

  The Style section of the Diagram Layout and Style Inspector contains controls for choosing and
  applying styles to a diagram on the canvas.

  \begin{description}[style=nextline]
    \item[{Choose\ldots{}}] Opens the diagram style chooser for applying a consistent style to all
      objects in the diagram. When you choose a style, it appears to the left in the
      \textbf{Style:} slot.
    \item[{Enforce Style}] Turn on to automatically apply the selected style to any new objects
      added to the canvas. (The styles of these objects can be subsequently changed.) With
      \textbf{Enforce Style} turned off, new objects added to the canvas retain their original
      style properties.
    \item[{Restyle Now}] Use to apply the diagram style to any objects that differ from it. This
      reverts any customizations you have made to object styles back to the default for that
      template. Changes you make to object styles are otherwise preserved regardless of how you
      reposition them in the diagram.
  \end{description}

  \begin{ogtip}
    Hold down the \textbf{Option} key to change Restyle Now to \textbf{Restyle Selection}, for more
    fine-grained control over the objects that receive the diagram style update.
  \end{ogtip}

  \subsection{Defining Your Canvases in the Canvas Data Inspector
  (\texorpdfstring{\textbf{PRO}}{PRO})}

  \begin{ogpro}
    You can also edit the note and other metadata of a canvas itself. Select the canvas in the
    Sidebar and then use the Canvas Data inspector to add metadata for the canvas.

    \ogfigure{0.800}{res/og7_inspectors62.png}

    As with the data added using the Object Data inspector, this information is used primarily for
    indexing and doesn't affect the appearance of your OmniGraffle document in any way.
  \end{ogpro}

  \section{The Document Inspectors}

  Use the Document inspectors to manage the properties of your OmniGraffle document.

  \subsection{The Margins Inspector}

  Here you can specify your own page margins, or choose \textbf{Use printer margins} to default to
  the margins defined by your printer driver (or by any custom settings you've made in
  \textbf{File} $\blacktriangleright$ \textbf{Page Setup}).

  \ogfigure{0.800}{res/og7_inspectors64.png}

  \subsection{Storing Information in the Document Data Inspector}

  The Document Data inspector has fields for lots of information about your document, in case you
  care to keep track of such things. The available fields are Subject, Copyright, Version,
  Description, and Comments. The pop-up menu includes options for adding information about the
  document's Authors, Organizations, Languages, Keywords, and Projects.

  \ogfigure{0.800}{res/og7_inspectors65.png}

  All of this data is made available to macOS's Spotlight searching feature, to help you find the
  diagram you're looking for.
  %%%%%%%%%%%%%%%%%%%%%%%%%%%%%%%%%%%%%%%%%%%%%%%%%%%%%%%%%%%%%%%%%%%%%%%%%%%%%%%%%%%%%%%%%%%%%%%%%%%
  \end{document}
